\chapter*{Preface to Volume II}
\addcontentsline{toc}{chapter}{Preface to Volume II}

Volume I treated CNA as an API and as eleven implementations of a rendering pipeline behind
that API. Volume II treats CNA as a \emph{project} --- a body of engineering practice, a set
of sibling libraries each with their own reasons to exist, and a working method for porting
real, historical games onto the result.

Part~I finishes the XNA namespace tour that Volume~I began with Graphics: input, audio,
device sensors, gamer services, networking, and the avatar system. Part~II turns to
\texttt{sharp-runtime} on its own terms, as the foundation every XNA-facing type in Volume~I
quietly depends on. Part~III does the same for \texttt{easy-gl} and \texttt{free-direct} ---
two libraries with almost opposite design philosophies (one deliberately general, one
deliberately narrow) that happen to back two of CNA's eleven graphics backends. Part~IV is
the cross-platform engineering that makes a single test suite mean something on Linux,
Windows-via-Wine, a browser tab, and an Android emulator at once, and the verification
discipline --- differential testing, pixel oracles, compile-time purity checks --- introduced
in Volume~I's first chapter, covered here in full mechanical detail. Part~V closes with the
practical work this whole apparatus exists to enable: migrating a real XNA/FNA game, a worked
case study of two actual legacy games (\emph{Speedy Blupi} and \emph{Planet Blupi}) ported via
\texttt{free-direct}, auditing compatibility claims against \texttt{xna4-spec} instead of
memory, and reading the project's own task-tracking conventions well enough to contribute to
it.

The same methodological commitment from Volume~I's preface holds here: a claim in this
volume is stated as fact because the project's own source, tests, or documentation state it
as a verified fact, with a mechanism behind it, not because it sounded right.
